\documentclass[../main.tex]{subfiles}
\begin{document}

Энэ бүлэгт системийн гол бүрэлдэхүүн хэсгүүдийг хэрхэн хэрэгжүүлснийг серверийн
талаас эхлэн клиент болон байршуулалтын түвшин хүртэл дэс дараатай тайлбарлав.
Бүх сервер талын код TypeScript хэл дээр, ESM модулийн стандартаар бичигдсэн
бөгөөд Fastify фреймворк дээр ажиллана.

\section{Баримт оруулах хоолой (Ingestion pipeline)}

Систем баримтыг хоёр замаар хүлээн авна. Эхнийх нь \texttt{POST /documents/upload}
буюу PDF файл байршуулах зам бөгөөд файлын дээд хэмжээг 20~MB-ээр хязгаарлаж,
mimetype эсвэл өргөтгөлөөр нь PDF мөн эсэхийг шалгана. PDF-ээс текст ялгахад
\texttt{pdf-parse} сан ашигласан бөгөөд уг сан нь CommonJS хэлбэрээр
нийтлэгдсэн тул ESM орчноос \texttt{createRequire} аргаар дуудаж холбов. Хэрэв
PDF дотроос текст олдохгүй бол (жишээ нь зураг хэлбэрээр хадгалагдсан скан
баримт) хэрэглэгчид ойлгомжтой алдааны мэдэгдэл буцаана. Хоёр дахь зам нь
\texttt{POST /documents/paste} буюу текст шууд буулгах зам бөгөөд 10-аас
500{,}000 тэмдэгтийн хязгаартай.

Баримт боловсруулах ажил нь Gemini API-ийн үнэгүй түвшний хязгаарлалтаас
шалтгаалан хэдэн минут үргэлжлэх боломжтой тул HTTP хүсэлтийн хариуг хүлээлгэж
болохгүй. Иймд хоёр зам хоёулаа баримтыг \texttt{pending} төлөвтэйгөөр
өгөгдлийн санд бүртгэж, ялгасан текстийг түр санах ойд хадгалаад, Node.js-ийн
\texttt{setImmediate} ашиглан боловсруулалтыг арын процесс болгон эхлүүлж,
клиентэд шууд 201 хариу буцаадаг болгов. Клиент талаас 3 секунд тутамд
\texttt{GET /documents/:id/status} хаягаар төлөвийг шалгаж, \texttt{ready}
болмогц баримтыг чатад ашиглах боломжтой болгоно. Бүрэн урсгалыг
Зураг~\ref{fig:ingestion-flow}-д үзүүлэв.

\begin{figure}[H]
  \centering
  \includegraphics[width=0.82\textwidth]{ingestion-flow.png}
  \caption{Баримт оруулах хоолойн урсгал ба төлөвийн шилжилт}
  \label{fig:ingestion-flow}
\end{figure}

Боловсруулалтын эхний алхам болох chunking буюу баримтыг жижиг хэсгүүдэд хуваах
ажлыг \texttt{chunker.ts} модуль гүйцэтгэнэ. Хэсэг бүрийн зорилтот хэмжээ 500
токен бөгөөд дараалсан хэсгүүд хоорондоо 15\%-ийн давхцалтай байхаар
тохируулав. Давхцал нь өгүүлбэрийн утга хэсгийн зааг дээр тасрахаас сэргийлж,
хайлтын чанарыг сайжруулдаг. Токены тоог үнэн зөв тооцохын тулд кирилл үсэгт
зориулсан тусгай тооцоо хэрэгжүүлэв: Gemini-ийн tokenizer кирилл текстийг латин
текстээс олон токенд хуваадаг тул текстийн 50\%-иас дээш хувь нь кирилл байвал
нэг токенд 2.5 тэмдэгт, бусад тохиолдолд 4 тэмдэгт гэж тооцно
(Код~\ref{lst:chunker}). Мөн хэсгийн таслах цэгийг сонгохдоо өгүүлбэрийн
төгсгөл, догол мөр, шинэ мөр, хоосон зай гэсэн дарааллаар хамгийн зохимжтой
байрлалыг хайдаг.

\begin{lstlisting}[caption={Кирилл-мэдрэг токен тооцоо (chunker.ts)}, label={lst:chunker}]
const TARGET_TOKENS = 500;
const OVERLAP_RATIO = 0.15;

function detectCharsPerToken(text: string): number {
  const cyrillicCount = (text.match(/[\u0400-\u04FF]/g) ?? []).length;
  return cyrillicCount / text.length > 0.5 ? 2.5 : 4;
}

export function estimateTokens(text: string): number {
  if (text.length === 0) return 0;
  return Math.ceil(text.length / detectCharsPerToken(text));
}
\end{lstlisting}

Хоёр дахь алхамд \texttt{embedder.ts} модуль хэсэг бүрийг Gemini-ийн
\texttt{gemini-embedding-001} загвараар 768 хэмжээст embedding вектор болгон
хөрвүүлнэ. Үнэгүй түвшний хүсэлтийн хязгаарлалтад (rate limit) тааруулан
хүсэлтүүдийг багцалдаг: нэг багцад дээд тал нь 15 хэсэг буюу ойролцоогоор 6000
токен багтааж, багц хооронд 3 секундын завсарлага авч, хүсэлт бүрд 30 секундын
timeout тавьсан. Хэрэв сервер 429 (хэт олон хүсэлт) эсвэл 5xx алдаа буцаавал
exponential backoff аргаар хамгийн ихдээ 8 удаа дахин оролддог бөгөөд Gemini-ийн
буцаасан \texttt{retryDelay} утгыг дагана. Энэ шийдэлд хүргэсэн асуудлын
дэлгэрэнгүйг \ref{sec:problem-large-paste} хэсэгт тайлбарлав.

Сүүлийн алхамд \texttt{vector-store.ts} модуль бүх хэсгийг embedding-ийн хамт
PostgreSQL-д хадгална. Хадгалалт нэг транзакц дотор 50 мөрийн багцаар хийгддэг
тул дундуур нь алдаа гарвал бүхэлдээ буцаагдаж (\texttt{ROLLBACK}), баримт
хагас дутуу хадгалагдах боломжгүй. Мөн оруулахын өмнө тухайн баримтын хуучин
хэсгүүдийг устгадаг тул давтан боловсруулалт аюулгүй (idempotent) юм.

\section{Хайлт ба хариу үүсгэлт}

Хэрэглэгчийн асуултад хариулахдаа систем эхлээд асуултыг мөн адил
\texttt{gemini-embedding-001} загвараар вектор болгож, дараа нь өгөгдлийн сангийн
\texttt{match\_chunks} функцээр хамгийн төстэй хэсгүүдийг хайна
(Код~\ref{lst:matchchunks}). Төстэй байдлыг косинус хэмжигдэхүүнээр тооцох
бөгөөд pgvector-ийн \texttt{<=>} оператор косинус зайг буцаадаг тул
$similarity = 1 - distance$ гэж хөрвүүлнэ. Хайлтыг хурдасгахын тулд
\texttt{chunks.embedding} багана дээр HNSW (Hierarchical Navigable Small World)
индекс үүсгэсэн. Функц нь \texttt{user\_id} болон сонгосон баримтуудын ID-аар
заавал шүүдэг тул нэг хэрэглэгчийн хайлт нөгөө хэрэглэгчийн баримтад хэзээ ч
хүрэхгүй.

\begin{lstlisting}[language=SQL, caption={match\_chunks функцийн гол хэсэг (001\_schema.sql)}, label={lst:matchchunks}]
SELECT id, document_id, content, page, chunk_index,
  1 - (embedding <=> query_embedding) AS similarity
FROM chunks
WHERE user_id      = match_user_id
  AND document_id  = ANY(match_document_ids)
  AND 1 - (embedding <=> query_embedding) > similarity_threshold
ORDER BY embedding <=> query_embedding
LIMIT match_count;
\end{lstlisting}

Хайлтын үр дүнг шууд ашиглахын оронд MMR (Maximal Marginal Relevance) буюу олон
талт байдлын дахин эрэмбэлэлт хэрэгжүүлэв. Эхлээд \texttt{match\_chunks}-ээс
дор хаяж 20 кандидат хэсэг авч, дараа нь greedy аргаар эцсийн $k=5$ хэсгийг
сонгоно. Кандидат $c$-ийн оноог

\[
score(c) = \lambda \cdot sim(c, q) - (1 - \lambda) \cdot \max_{s \in S} J(c, s)
\]

томьёогоор тооцно. Энд $sim(c,q)$ нь асуулттай харьцуулсан косинус төстэй
байдал, $J(c,s)$ нь аль хэдийн сонгогдсон $S$ олонлогийн хэсгүүдтэй
харьцуулсан Jaccard төстэй байдал (үгийн олонлогийн огтлолцлын харьцаа),
$\lambda$ нь хамаарал ба олон талт байдлын тэнцвэрийг тохируулах параметр юм.
Ингэснээр бараг ижил агуулгатай хэсгүүд давхардан сонгогдохоос сэргийлж,
хариултын context олон талын мэдээлэл агуулдаг болно. Хэрэглэгч MMR-ийг
унтраавал (\texttt{useMMR=false}) энгийн similarity эрэмбийн эхний $k$ хэсгийг
буцаана. Босго болон $\lambda$-гийн анхдагч утгыг $0.1$ ба $0.5$ гэж
тохируулсан шалтгааныг \ref{sec:problem-crosslingual} хэсэгт тайлбарлав.

Хариу үүсгэлтийг \texttt{generator.ts} модуль гүйцэтгэнэ. Олдсон хэсгүүдийн
агуулгыг \texttt{---} тусгаарлагчаар нэгтгэн context болгож, дараах системийн
зааврын хамт Gemini рүү илгээнэ:

\begin{quote}
\emph{«Та баримтын туслах. Хэрэглэгчийн асуултад зөвхөн өгөгдсөн баримтын
context ашиглан хариул. Хариулт бичихдээ баримтын хэл болон бичгийн системийг
дагаж бич --- баримт кирилл монгол бол кирилл монголоор, латин бол бас кирилл
монголоор хариул. [\ldots] Асуулт ямар ч хэлээр байсан хамаагүй. Хариулт
context дотор байхгүй бол тодорхой хэл. Товч, ойлгомжтой бай.»}
\end{quote}

Заавар нь загварыг зөвхөн context-оос хариулахыг шаардаж hallucination буюу
байхгүй мэдээлэл зохиохоос сэргийлэхийн зэрэгцээ хариултын хэлийг асуултын
хэлээс бус \emph{баримтын} хэлнээс хамааруулдаг нь монгол баримттай ажиллах
үеийн чухал дүрэм болсон. Найдвартай байдлын үүднээс загваруудыг
\texttt{gemini-2.5-pro} $\rightarrow$ \texttt{gemini-2.5-flash} $\rightarrow$
\texttt{gemini-2.5-flash-lite} гэсэн fallback гинжээр дуудна: аль нэг загвар
алдаа буцаавал дараагийнх нь оролдоно, мөн загвар бүр 429 эсвэл 503 алдааны
үед 2 секунд хүлээгээд нэг удаа дахин оролддог. Хэрэв хайлтаас тохирох хэсэг
огт олдоогүй бол Gemini-г дуудалгүйгээр \emph{«Сонгосон баримтуудаас энэ
асуултад хамаарах хангалттай мэдээлэл олдсонгүй»} гэсэн шударга хариуг буцаана.

Хариулт бүрийг эх сурвалжийн хамт хадгална: \texttt{messages.sources} JSONB
баганад \texttt{\{chunkId, documentId, content, page, chunkIndex,
similarity\}} бүтэцтэй массив бичигдэж, клиент талд хариултын доор эх
сурвалжийн картууд болон харагдана. Чатын бүрэн урсгалыг
Зураг~\ref{fig:chat-sequence}-д үзүүлэв.

\begin{figure}[H]
  \centering
  \includegraphics[width=\textwidth]{chat-sequence.png}
  \caption{Чатын хүсэлтийн бүрэн дараалал (sequence diagram)}
  \label{fig:chat-sequence}
\end{figure}

\section{Нэвтрэлт ба олон хэрэглэгчийн тусгаарлалт}

Нэвтрэлтийн системийг Better Auth сангаар хэрэгжүүлэв. Better Auth нь тусдаа
үйлчилгээ биш, API процессын дотор ажиллаж, аппликейшнтэйгээ нэг PostgreSQL
Pool-ийг хуваалцан ашигладаг. Fastify-ийн хүсэлтийг Fetch API хэлбэрт хөрвүүлэн
\texttt{/api/auth/*} замуудад холбосон. Хэрэглэгчийн бүртгэл имэйл + нууц үг
хослолоор хийгдэх бөгөөд нууц үг нь scrypt алгоритмаар hash хийгдэж хадгалагдана
(хамгийн багадаа 8 тэмдэгт). Better Auth-ийн үүсгэдэг ID-г
\texttt{advanced.database.generateId} тохиргоогоор UUID болгосноор аппликейшны
хүснэгтүүдийн \texttt{user\_id UUID} баганатай шууд нийцдэг.

Session удирдлага клиент бүрд өөр өөрөөр шийдэгдэв: вэб клиент session
cookie-г хөтчийн стандарт механизмаар (\texttt{credentials: "include"})
илгээдэг бол мобайл клиент cookie-г Expo SecureStore-д шифрлэн хадгалж,
\texttt{@better-auth/expo} plugin-ий \texttt{authClient.getCookie()} функцээр
хүсэлт бүрийн Cookie толгойд хавсаргана.

Хамгаалалтын гол зарчим нь давхар шүүлт юм. Нэгд, хамгаалагдсан бүх route
эхэндээ \texttt{requireUser()} функцээр session-ийг шалгаж, нэвтрээгүй
хэрэглэгчид 401 буцаана (нээлттэй үлдсэн нь зөвхөн \texttt{/health},
\texttt{/config}, \texttt{/api/auth/*}). Хоёрт, өгөгдлийн сангийн \emph{бүх}
асуулга --- баримтын жагсаалтаас эхлээд \texttt{match\_chunks} функцийн дотоод
шүүлт хүртэл --- \texttt{user\_id}-аар хязгаарлагддаг. Ингэснээр нэг
хэрэглэгчийн асуулт, хайлт, түүх нөгөө хэрэглэгчийн өгөгдөлд ямар ч нөхцөлд
хүрч чадахгүй. Бүртгэлээс хамгаалагдсан хүсэлт хүртэлх бүрэн дарааллыг
Зураг~\ref{fig:auth-sequence}-д үзүүлэв.

\begin{figure}[H]
  \centering
  \includegraphics[width=\textwidth]{auth-sequence.png}
  \caption{Нэвтрэлтийн дараалал: бүртгэл, session үүсэлт, 401 хамгаалалт}
  \label{fig:auth-sequence}
\end{figure}

\section{Вэб клиент}

Вэб клиентийг React, Vite, Tailwind CSS ашиглан хөгжүүлэв. Апп нь танилцуулга
(landing) хуудас болон ажлын талбар (workspace) гэсэн хоёр үндсэн хэсэгтэй,
чиглүүлэлтийг React Router гүйцэтгэнэ. Хэрэглэгчид харагдах бүх текстийг монгол
хэлээр бичиж, кирилл үсгийг зөв дүрслэхийн тулд фонтын жагсаалтад Noto Sans-ыг
нэмсэн. Дизайн нь glassmorphism хэв маягийн нэгдсэн токенуудад
(өнгө, радиус, зай) суурилдаг.

API-ийн хаягийг тогтмол бичихийн оронд ажиллаж буй орчноос нь тодорхойлдог
болгов: \texttt{VITE\_API\_BASE\_URL} орчны хувьсагч байхгүй үед
\texttt{window.location.hostname} дээр 4000 портыг залгаж ашиглана. Ингэснээр
нэг дотоод сүлжээний өөр төхөөрөмжөөс (жишээ нь утаснаас) вэб аппыг нээхэд API
дуудлага автоматаар зөв хост руу очдог. Баримтын төлөвийг 3 секунд тутамд
шалдаг polling нь зөвхөн \texttt{pending} эсвэл \texttt{processing} төлөвтэй
баримт байгаа үед л ажиллаж, бүгд \texttt{ready} болмогц зогсдог.

Чатын оролтын хэсэгт (Composer) хайлтын тохиргоог хэрэглэгчид ил гаргаж өгөв:
«Олон талт хариу» товч нь MMR-ийг асааж унтраадаг бол тохиргооны цонхонд
«Хамаарлын босго» (threshold) болон «Олон талт байдал ($\lambda$)» гэсэн хоёр
гулсагч 0--1 хооронд 0.05 алхамтайгаар утга сонгуулж, анхдагч утга руу буцаах
товчтой. Эдгээр утга \texttt{/chat} хүсэлтийн бие дотор дамждаг.

% TODO(Тэнгис): вэб дэлгэцийн зургуудыг figures/ дотор хийгээд доорх мөрүүдийг тайлна уу.
% \begin{figure}[H]
%   \centering
%   \includegraphics[width=\textwidth]{web-workspace.png}
%   \caption{Ажлын талбар: чат, баримтын жагсаалт, эх сурвалжийн картууд}
%   \label{fig:web-workspace}
% \end{figure}
%
% \begin{figure}[H]
%   \centering
%   \includegraphics[width=0.7\textwidth]{web-settings.png}
%   \caption{Хайлтын тохиргоо: MMR товч ба threshold/$\lambda$ гулсагчид}
%   \label{fig:web-settings}
% \end{figure}

\section{Мобайл клиент}

Мобайл клиентийг Expo дээр React Native ашиглан хөгжүүлэв. Дэлгэцүүдийн
чиглүүлэлтийг файлд суурилсан Expo Router, загварчлалыг NativeWind (React
Native-д зориулсан Tailwind) гүйцэтгэнэ. Апп нээгдэхэд танилцуулга
(onboarding) $\rightarrow$ нэвтрэх $\rightarrow$ ажлын талбар гэсэн дарааллаар
чиглүүлдэг бөгөөд танилцуулгыг алгасах сонголт зөвхөн хүчинтэй session байгаа
үед үйлчилнэ. Чиглүүлэлтийн бүрэн бодлогыг Зураг~\ref{fig:mobile-routing}-д
үзүүлэв; танилцуулгаас эхлэлийн \texttt{"/"} зам руу буцах шилжилт огт
байхгүй байх нь \ref{sec:problem-bootloop} хэсэгт өгүүлэх boot loop алдааны
засварын гол зарчим юм.

\begin{figure}[H]
  \centering
  \includegraphics[width=0.9\textwidth]{mobile-routing.png}
  \caption{Мобайл аппын чиглүүлэлтийн төлөвийн диаграм}
  \label{fig:mobile-routing}
\end{figure}

Мобайл орчны нэг онцлог бэрхшээл нь API серверийн хаяг тодорхойгүй байдаг явдал
юм: Android эмулятор хост машиныг \texttt{10.0.2.2} гэж хардаг бол бодит утас
дотоод сүлжээний IP хаяг шаарддаг. Үүнийг шийдэхийн тулд \texttt{api-base.ts}
модуль боломжит бүх хаягийг зэрэг шалгадаг: \texttt{EXPO\_PUBLIC\_API\_URLS}
жагсаалт, Metro bundler-ийн хостын IP, эмуляторын хаяг, localhost гэсэн
кандидатуудад 2 секундын timeout-той \texttt{GET /health} хүсэлтийг зэрэг
илгээж, хамгийн түрүүнд амьд хариу өгсөн хаягийг сонгоно
(Код~\ref{lst:apibase}). Сонгогдсон хаяг session-ий туршид кэшлэгдэх ба
сүлжээний алдаа гарвал кэш цэвэрлэгдэж дахин шалгалт хийгддэг.

\begin{lstlisting}[caption={Зэрэгцээ liveness шалгалт (api-base.ts)}, label={lst:apibase}]
resolving = Promise.any(candidateList().map(probe))
  .then((base) => {
    resolvedBase = base;
    return base;
  })
  .finally(() => {
    resolving = null;
  });
\end{lstlisting}

Чатын дэлгэц, баримт сонгох цонх, харилцан ярианы түүхийн хажуугийн цэс
(зүүн ирмэгээс чирч нээгддэг) зэрэг нь вэб хувилбартай ижил үйлдлүүдийг
гүйцэтгэдэг бөгөөд Composer-ийн «Олон талт хариу» товч, threshold/$\lambda$
гулсагч бүхий тохиргооны цонх мөн бүрэн ижил ажиллана.

% TODO(Тэнгис): мобайл дэлгэцийн зургуудыг (Expo Go эсвэл эмулятороос) авч
% figures/ дотор хийгээд доорх мөрүүдийг тайлна уу.
% \begin{figure}[H]
%   \centering
%   \includegraphics[width=0.32\textwidth]{mobile-onboarding.png}\hfill
%   \includegraphics[width=0.32\textwidth]{mobile-chat.png}\hfill
%   \includegraphics[width=0.32\textwidth]{mobile-sidebar.png}
%   \caption{Мобайл апп: танилцуулга, чат, хажуугийн цэс}
%   \label{fig:mobile-screens}
% \end{figure}

\section{Docker байршуулалт}
\label{sec:impl-docker}

Хөгжүүлэлтийн болон байршуулалтын орчныг Docker Compose-оор нэгтгэв. Стек нь
гурван үндсэн үйлчилгээтэй: \texttt{postgres}
(\texttt{pgvector/pgvector:pg16} image, зөвхөн \texttt{127.0.0.1}-д
холбогдсон), \texttt{api} (4000 порт) болон \texttt{web} (5173 порт). Мөн
өгөгдлийн санг вэбээс удирдах pgAdmin-ийг сонголтоор (\texttt{--profile
tools}) асаах боломжтой. Өгөгдөл named volume-д хадгалагддаг тул image дахин
build хийхэд устдаггүй. Бүрэн байршуулалтын топологийг ---
хэрэглэгчийн төхөөрөмжүүд, контейнерууд, volume-ууд болон гадаад Gemini
API-ийн холболтыг --- Зураг~\ref{fig:deployment}-д үзүүлэв.

\begin{figure}[H]
  \centering
  \includegraphics[width=\textwidth]{deployment.png}
  \caption{Байршуулалтын диаграм: Docker Compose орчин}
  \label{fig:deployment}
\end{figure}

Image build хийх хугацааг багасгахын тулд Docker-ийн layer cache-д тааруулсан
дараалал ашиглав: эхлээд зөвхөн багц менежерийн файлуудыг хуулж хамаарлуудыг
суулгаад, кодыг хамгийн сүүлд хуулна (Код~\ref{lst:dockerfile}). Ингэснээр код
өөрчлөгдөх бүрд \texttt{pnpm install} алхам cache-ээс шууд авагдана. Анхаарах
нэг зүйл нь monorepo-ийн \emph{бүх} workspace-ийн \texttt{package.json}
файлыг энэ давхаргад оруулах ёстой: нэгийг нь орхиход lockfile тохирохгүй
болж cache байнга эвдэрдэг байсныг илрүүлж засав.

\begin{lstlisting}[caption={Layer cache-д зориулсан Dockerfile бүтэц (apps/api/Dockerfile)}, label={lst:dockerfile}]
COPY package.json pnpm-lock.yaml pnpm-workspace.yaml turbo.json .npmrc ./
COPY apps/web/package.json apps/web/package.json
COPY apps/api/package.json apps/api/package.json
COPY apps/mobile/package.json apps/mobile/package.json
COPY packages/shared/package.json packages/shared/package.json
COPY packages/api-client/package.json packages/api-client/package.json
RUN pnpm install --frozen-lockfile
COPY . .
\end{lstlisting}

Өгөгдлийн сангийн бүтцийн өөрчлөлтийг удирдахад migration систем хэрэгжүүлэв.
Анхны бүтэц \texttt{infra/postgres/init/} доторх SQL файлуудаар зөвхөн шинэ
volume дээр нэг удаа үүсдэг бол дараагийн бүх өөрчлөлт
\texttt{infra/postgres/migrations/} доторх дугаарлагдсан
(\texttt{NNN\_нэр.sql}) файлууд болж нэмэгдэнэ. API асах бүрдээ
\texttt{schema\_migrations} хүснэгттэй тулгаж, хэрэгжээгүй migration бүрийг
тус тусын транзакцаар нэг л удаа ажиллуулна. Энэ нь production орчинд өгөгдөл
устгалгүйгээр бүтцийн өөрчлөлт хийх боломжийг олгодог.

\end{document}
